\chapter{Introducción}

La creciente digitalización de las operaciones institucionales ha ampliado de forma sostenida la superficie de exposición de las organizaciones frente a ataques informáticos. Dentro de este panorama, las Amenazas Persistentes Avanzadas (APT) ocupan un lugar particular: se trata de campañas dirigidas, ejecutadas por adversarios con recursos y objetivos definidos, que se despliegan de manera sigilosa a lo largo de meses y en múltiples etapas sucesivas \parencite{hutchins2011killchain}. Su impacto se extiende más allá del incidente puntual y compromete de forma prolongada la confidencialidad de la información de la organización afectada.

La dificultad para enfrentarlas radica en que los mecanismos de detección predominantes fueron diseñados para identificar eventos aislados y no secuencias de acciones distribuidas en el tiempo, razón por la cual las APT logran evadirlos con frecuencia \parencite{alshamrani2019survey}. Esta limitación ha motivado el interés por enfoques capaces de modelar la progresión de un ataque como un proceso, incorporando la incertidumbre inherente al comportamiento del adversario. Entre ellos, los modelos probabilísticos ofrecen un marco formal para estimar en qué etapa de una campaña se encuentra un atacante a partir de la evidencia observable en la red.

El presente proyecto aborda esta problemática mediante la implementación y validación de un modelo basado en Cadenas Ocultas de Markov (HMM) para la detección de APT, evaluado empíricamente sobre dos conjuntos de datos complementarios: DAPT2020 \parencite{myneni2020dapt2020}, que permite la comparación con los enfoques reportados en la literatura, y Unraveled \parencite{myneni2023unraveled}, que reproduce condiciones de operación más realistas mediante la emulación de seis semanas de actividad organizacional. Ambos fueron seleccionados mediante criterios explícitos derivados de los requerimientos del modelo: campañas APT multietapa, etiquetas de etapa por registro, tráfico de red a nivel de flujo y actividad benigna continua. La investigación se enmarca en la metodología de Design Science Research \parencite{peffers2007dsrm}, orientada a la construcción de artefactos que resuelven problemas concretos, y se apoya en una revisión sistemática de la literatura para fundamentar las decisiones de diseño del modelo propuesto.

Este documento corresponde al Informe Final de Seminario de Título y se organiza de la siguiente manera. El Capítulo~2 describe el problema abordado y su relevancia. El Capítulo~3 define el objetivo general y los objetivos específicos del proyecto. El Capítulo~4 presenta el marco teórico con los conceptos y técnicas necesarios para comprender la propuesta. El Capítulo~5 detalla la metodología de investigación adoptada y el procedimiento seguido en la revisión sistemática de la literatura. El Capítulo~6 expone el estado del arte resultante de dicha revisión. El Capítulo~7 presenta el plan de trabajo y su cronograma para ambos semestres. El Capítulo~8 describe la propuesta de solución y su diseño experimental. El Capítulo~9 presenta una primera implementación de la propuesta y sus resultados preliminares sobre ambos conjuntos de datos. El Capítulo~10 recoge las conclusiones de esta etapa y el Capítulo~11 declara el uso de herramientas de inteligencia artificial en la elaboración del documento.